\chapter{Summary}\label{ch:work-brief}

This document describes the Apple M5 GPU (the G17) at the level beneath Metal: its instruction bytes, what
they do, and how work reaches the hardware. From that description the project built a compiler that emits native G17
code, and two ways to run it: through Metal, and below Metal, where the project's own runtime submits work through Apple's
private IOGPU interfaces and Metal is never loaded (\cref{fig:bypass}).

\section{Measured hardware and software}\label{sec:part-software}

Every measurement here was made on one machine: an Apple M5 Pro (the H17s GPU, applegpu\_g17s) with 20 GPU cores and
48~GB of memory. The stated platform is macOS 26.6.2 (build 25G83) with the Xcode 27.0 toolchain (Metal front end
32023.921, against an OS compiler at 32023.886; MM~\mm{22}, MM~\mm{25.91}). The below-Metal records carry that OS build in their
own evidence. Most Metal-path evidence files record no OS build, so for those results the build is the stated
platform, not a recorded one.

A result on this part does not establish behaviour on another G17 product, another OS build or another generation
(\cref{sec:measured-part-software}). Where something has been checked across parts or builds, the text says so.

\section{Recovered behaviour}\label{sec:has-been-recovered}

Apple documents how to program the GPU and its neural accelerators through Metal, and nothing beneath that interface.
The project has recovered much of that lower layer and tested it on the hardware.

\begin{itemize}
\item In the register file, the usable namespace ends at R125, and an instruction naming any register from R126 to R143
  is dropped whole, without an error. Register operands carry a lifetime modifier, which a few short forms fix. A \emph{release} frees the register as it is
  read, and a later reader must not rely on its contents: controlled probes read zero, but what a released register
  returns is not established in general. These rules constrain register allocation, loops and composition (\cref{ch:register-file}).
\item Instructions are 2 to 22 bytes in even steps, and the opcode does not fix the length: each instance
  carries its own. Apple's code uses every even length from 2 to 16, plus 18- and 22-byte forms such as \op{10909} and \op{14661}; no 20-byte instruction has been seen. Operands are linear functions of
  scattered bits, recovered as field maps. The tensor-instruction encoder reproduces 800 of 800 compiler-emitted
  instructions byte for byte. Of the \val{named-instructions} opcodes the project names, \val{executed-instructions} are confirmed by execution
  (\cref{ch:instruction-encoding}).
\item For the tensor unit, Apple's decoder admits ten tensor MMA forms. The mapping from each lane's registers to matrix elements
  was measured with one-hot probes over 6,651 dispatches. For half inputs, the exact reduction order of one 16 × 16 × 16
  multiply-accumulate, including where the accumulator enters, reproduces 5,120 of 5,120 tile elements bit for bit,
  where a sequential chain, a balanced tree and an exactly rounded sum each match fewer than 2,800; other input types are
  recorded at single points (\cref{ch:tensor-unit}).
\item Ordering uses no wait instruction. The producer of a late value names one of eight scoreboard slots, and
  the consumer carries an eight-bit wait mask. If the slot's bit is left out, the consumer reads the register before the value lands
  (zeros, in the measured tensor case) and nothing reports an error (\cref{ch:synchronization}).
\item Conversions, fused operations and special-function instructions have been measured, some of them input by
  input. The same chapters give the address law, the load and store forms, threadgroup memory, the imageblock,
  constants and atomics (\crefrange{ch:scalar-arithmetic}{ch:synchronization}).
\item In the kernel image, every byte across 421 delivered archives has a known source. Apple's loader requires 148
  constants, and the meaning of 12 of their bytes is still open. Metal checks the container's format and does not inspect the code in it (\cref{ch:kernel-metadata-image}).
\item In the submission machinery, the launch words, the Submit record, the kernel entry and the completion path have
  been recovered well enough to run complete computations without Metal, within the measured resource graphs
  (\cref{ch:submission-below-metal}).
\end{itemize}

\section{Project components and execution paths}\label{sec:project-authors-two}

The project authors the IR, the compiler (instruction selection, register allocation, encoding and tensor lowering),
the native instruction bytes, the object, metadata and container around them, an executor for each path, and the
verification tools. \Cref{fig:bypass} shows both paths and names the Apple components each one meets.

% Figure 0.1: the two execution paths. Blue: the project's components. Red: the
% project's native runtime. Grey: Apple's components.
\begin{figure}[htbp]
\centering
\begin{tikzpicture}[node distance=6mm and 8mm]
  \node[ours, minimum width=46mm] (ir) {The project's IR and kernel builders};
  \node[ours, minimum width=46mm, below=of ir] (cc) {The project's compiler};
  \node[ours, minimum width=46mm, below=of cc] (code) {The project's native G17 instruction bytes};
  \draw[flow] (ir) -- (cc);
  \draw[flow] (cc) -- (code);

  \node[ours, text width=52mm, below left=9mm and -14mm of code] (img)
    {The project's \mbox{Metal-compatible} image\\[-1pt]{\footnotesize object, metadata, library, archive}};
  \node[ours, text width=52mm, below=of img] (exec)
    {The project's Metal-path executor\\[-1pt]{\footnotesize \code{decodegen}}};
  \node[apple, text width=52mm, below=of exec] (metal)
    {Apple Metal\\[-1pt]{\footnotesize pipeline creation and submission}};

  \node[nat, text width=52mm, below right=9mm and -14mm of code] (rt)
    {The project's native runtime\\[-1pt]{\footnotesize allocation, resource and launch state,}\\[-2pt]
     {\footnotesize submission records, completion handling}};

  \draw[flow] (code.south) -- ++(0,-4mm) -| (img.north);
  \draw[flow] (code.south) -- ++(0,-4mm) -| (rt.north);
  \draw[flow] (img) -- (exec);
  \draw[flow] (exec) -- (metal);

  \path (metal.south) ++(0,-9mm) coordinate (iogpuTop);
  \node[band, below=9mm of metal.south -| code, minimum width=118mm] (iogpu)
    {Apple private IOGPU framework and IOKit interfaces};
  \node[band, below=5mm of iogpu, minimum width=118mm] (drv) {Apple AGX kernel driver};
  \node[band, below=5mm of drv, minimum width=118mm] (fw) {Apple GPU firmware};
  \node[box, draw=ink, fill=white, below=5mm of fw, minimum width=118mm] (gpu)
    {M5 GPU: scalar and tensor execution};

  \draw[flow] (metal.south) -- (metal.south |- iogpu.north);
  \draw[->, line width=1.4pt, color=native] (rt.south) -- (rt.south |- iogpu.north)
    node[pos=0.13, right=1mm, tag, text=native, text width=26mm, align=left]
      {direct submission:\\bypasses Metal};
  \draw[flow] (iogpu) -- (drv);
  \draw[flow] (drv) -- (fw);
  \draw[flow] (fw) -- (gpu);

  % Path labels
  \node[font=\small\bfseries, text=ours, anchor=base west, inner sep=1pt] at ([yshift=2mm]img.north west) {Through Metal};
  \node[font=\small\bfseries, text=native, anchor=base east, inner sep=1pt] at ([yshift=2mm]rt.north east) {Below Metal};

  % The Metal layer, which the native path does not enter
  \begin{scope}[on background layer]
    \node[draw=apple, dashed, rounded corners=3pt, fill=apple!4, inner sep=3mm,
          fit=(metal) (metal -| rt.east)] (mlayer) {};
  \end{scope}
  \node[tag, text width=44mm] at ($(metal.east)!0.5!(metal.east -| rt.south)$) {Metal is not loaded in\\the native runtime's process};

  % Brace over the shared Apple layers
  \draw[decorate, decoration={brace, amplitude=5pt}, line width=0.6pt, color=apple]
    ([xshift=3mm]iogpu.north east) -- ([xshift=3mm]fw.south east)
    node[midway, right=6pt, tag, text width=18mm, align=left] {below both\\paths};
\end{tikzpicture}
\caption{The two execution paths. Both execute the project's own native GPU code. The native runtime bypasses Metal's
pipeline creation and submission, builds the measured resource and launch state itself, and submits through Apple's
private IOGPU interfaces. Apple's kernel driver and firmware remain in use. Blue: the project's components; red: the
project's native runtime; grey: Apple's.}
\label{fig:bypass}
\end{figure}


\begin{itemize}
\item Through Metal (\cref{ch:through-metal}), the project hands Metal an image. Metal builds the pipeline from its
  archived code and submits the work.
\item Below Metal (\cref{ch:below-metal-native}), the native runtime takes the program bytes directly. It uploads them,
  writes each launch's resource and launch state, builds the Submit records, and calls IOGPU itself. No Metal
  pipeline, command buffer or device exists in that process.
\item Below both paths, the same Apple components remain: IOGPU, the AGX kernel driver and the firmware.
\end{itemize}

Build and validation use one more Apple component. The compiler's encoders, its whole-program
checks and the CPU emulator decode programs with Apple's G17 decoder, loaded from \code{GPUCompiler.framework} (
\cref{ch:instruction-encoding}). That decoder is not used when the GPU runs a program, but it is the reason compiling and
validating currently require macOS. Metal source also passes through Apple's \code{xcrun metal} front end before it reaches the
project's IR (\cref{ch:source-native-program}).

\section{The running example}\label{sec:one-program-ir}

The running example computes a 17 × 19 × 16 matrix product, half-precision inputs and fp32 output, and then adds 1.0
to one output element. Later chapters come back to it at each layer and check it there. The bytes are the project compiler's output. Recompiled on the CPU, the program is byte-identical to
the one that ran on the GPU below Metal: 1,232 bytes, SHA-256 \hash{8bb67a3c65919cb7fa9a6f2e70d152cb2d8086fba0cade7437f18d1d8114c667},
register count 76.

The source is seven lines of the project's IR: three buffers, \code{tensor\_matmul(A, B, C, M=17, N=19, K=16)}, a
load, an fp32 add and a store (\cref{ch:running-example-through}).

The 17 × 19 output does not divide into 16 × 16 tiles, so the product is two tile rows by two tile columns, and
the second row and column are partial. The compiler emits 101 instructions, grouped by phase in \cref{tab:example-phases}.

\begin{xltabular}{\linewidth}{@{}l>{\hsize=1.10\hsize}X>{\hsize=0.90\hsize}X@{}}
\caption{Instructions of the 17 × 19 × 16 example program by phase, as read with Apple's decoder, with what each group does.}\label{tab:example-phases}\\
\toprule
Phase & Instructions & What they do \\
\midrule
\endfirsthead
\tablecontinued{3}
\toprule
Phase & Instructions & What they do \\
\midrule
\endhead
addresses & 1 \op{14060}, of the lane index, then
shifts, ANDs and adds & compute each lane's row, column and byte offset (
\cref{sec:fragment-layouts}) \\*
masks & 20 \op{612} and \op{437} & enable only the 17 rows and 19 columns that exist (
\cref{sec:masks-predicated-memory}) \\*
loads & 4 \op{12674} and 4 \op{12675} & gather two A and two B fragments, eight bytes per lane each \\*
multiply & 4 \op{5107} & one per output tile, since K = 16 is one slice \\*
stores & 2 \op{17257} and 6 \op{17258} & write the four tiles, masked where partial \\*
epilogue & \op{14059}, \op{12682}, \op{11842}, \op{998}, \op{17229}, \op{684} & add 1.0 to one element \\*
\bottomrule
\end{xltabular}

The first \op{5107}, is ten bytes,
\hash{a70405112204a4620202}. Apple's decoder reads it as follows:

\begin{lstlisting}
op5107  R24..R31 <- 0x01000000, 1, R20_R21_R22_R23, 0, 2, R68_R69_R70_R71, 0, 2
        D        flags word        A fragment      lifetime, type   B fragment   lifetime, type
\end{lstlisting}

\begin{itemize}
\item \textbf{D, R24–R31,} is an eight-register fp32 accumulator group: the output tile at application rows 0–15,
  columns 0–15. In the instruction's canonical coordinates lane 5 holds elements (4, 4)–(4, 7) in its first four
  registers and (5, 4)–(5, 7) in the next four. The compiler packs every tile with the B map, so its store writes the
  same registers to application rows 2 and 10, columns 4–7. The two labelings differ by a one-bit rotation of the row
  index (\code{rotl1(2) = 4}, \code{rotl1(10) = 5}); A is packed the same way, so the rotation only relabels the free M index
  and the product is correct in application coordinates (\cref{sec:coordinate-systems}).
\item \textbf{A, R20–R23, and B, R68–R71,} are four-register fragments of 16-bit values. Type code 2 means half.
\item \textbf{The flags word 0x01000000} has bit 24 set: wait for scoreboard slot 0. All eight tensor loads write slot 0. Each
  load's own control word, 0x01100000, both publishes slot 0 and waits on it (\cref{sec:scoreboard-publish-wait,sec:scoreboard}). On an MMA
  whose loads fill slot 7, clearing the matching wait bit made the multiply read its fragments before the loads landed,
  and it produced zeros (measured, MM~\mm{6}); that control was run on slot 7, not on this program.
\item Each fragment is read by two of the four MMAs. Its first reader leaves it live (modifier 0)
  and its second, the last use, releases it (16); nothing reads the register after that (\cref{sec:release-does}).
\item \textbf{The 1 after the flags word} is a constant: no encoding bit varies it (\cref{sec:field-maps}).
\item Each output element sums 16 products in a fixed tree: adjacent pairs, then pairs eight apart,
  rounding to nearest-even fp32 at every step (the rule fitted on half inputs, which this program uses). This from-zero form starts from the first of the four partial sums,
  where the accumulating form adds C first (\cref{sec:arithmetic-semantics-supported}).
\end{itemize}

Read from the bytes, the epilogue is ordered by these waits:

\begin{enumerate}
\def\labelenumi{\arabic{enumi}.}
\item The \op{14059} of the threadgroup position publishes slot 0.
\item The \op{12682} waits on slot 0 for its index register and publishes slot 7.
\item The \op{998} waits on slot 7 for the loaded value.
\item The store needs no wait: its value comes from an ALU instruction, and ALU results are interlocked (\cref{sec:dependences-issue-scalar}).
\end{enumerate}

Alongside the code, the compiler reports what an executor needs: the three bindings with their element types, which
binding is written, the 24 distinct (opcode, length) forms used, an empty constant pool, the register count and
the system registers read. The image author turns that into the object, metadata, library and archive Metal loads (\cref{ch:kernel-metadata-image}).

The program was then run on the GPU below Metal (measured; \code{docs/g17-first-dispatch-trials.md}, "A second authored tensor program
below Metal"). The native runtime wrote the program at allocation 0, offset 0x6c0. Relative to the first tensor
program it had run, the measured resource graph needed only these field changes:

\begin{itemize}
\item two moved view offsets;
\item one C pointer;
\item the kernel word at offset 0x3f8, from 0xa to 0xc;
\item two binding pointers;
\item A, B and C windows of 544, 608 and 1,292 bytes, which are exactly 17 × 16 halves, 16 × 19 halves and 17 × 19 floats.
\end{itemize}

Three fresh processes each submitted the program through IOGPU with no Metal loaded. In each, 323 of 323 fp32 outputs
were bit-exact, inputs and guards were intact, both completion callbacks arrived and no recovery event occurred. The
outputs hash to \hash{ec43027260f5a0b8dd59526c17102ffe00066852f10af42a91103c36e778faf6}, the same as an independently
computed reference and as the earlier run through Metal (\cref{ch:submission-below-metal}). That is one program at one
shape; how far each layer generalises is stated in the chapters cited above.

\section{Applications}\label{sec:understanding-supports}

\begin{itemize}
\item InternLM2.5-1.8B decodes on the compiler's kernels, run through Metal. Single
  stream at a 1,024-token context, it runs at 1.14× the decode rate of mlx-lm's \code{generate\_step}. Several other models and workloads are
  measured, each with its own baseline and limits (\href{https://github.com/sbryngelson/AGXForge/blob/main/docs/showcase.md}{the showcase}).
\item The compiler was also run on code it was not written for. Of 125 compiled sources from a frozen 197-source
  census, 101 have run and been checked on the hardware, and none crashes the driver (MM~\mm{25.149.1}).
\item A CPU emulator, working from the program bytes, executes one position of a shipped batched-decode
  graph and of single-stream decode, dispatch by dispatch: 316 of 316 and 172 of 172 dispatches byte-identical to the
  GPU, and a 16~MB prefill slice identical (MM~\mm{25.197}). It is an internal verification tool, and it uses Apple's
  decoder.
\item MiniLM-L6 sentence retrieval and Qwen2.5-0.5B-Instruct generation ran through
  the native submission path, each with independently checked outputs (MM~\mm{25.210}; \cref{ch:below-metal-native}).
\end{itemize}

\section{Limits of the evidence}\label{sec:where-evidence-stops}

\begin{itemize}
\item Nothing here is claimed for other parts or builds, and the OS build is
  verified per record only where a record carries it (\cref{sec:part-software,sec:measured-part-software}).
\item The InternLM2 decode lead over mlx-lm (\cref{sec:understanding-supports})
  is due to the kernels' design. In a matched study, Apple's compiler, given the same kernels, ran them 10–15 percent
  faster with identical tokens, so the lead does not show that native instruction control is needed. The tensor-unit
  kernels were not part of that study (MM~\mm{25.211}).
\item Below Metal, complete computations run with checked results, but not faster
  than the compared Apple paths (MM~\mm{25.210.5}).
\item \Cref{ch:below-metal-native} keeps four kinds of evidence apart: stage checks
  against actual inputs, comparisons with the original checkpoint, exact differentials, and one auxiliary control,
  which fails.
\item Validation uses Apple's decoder. An independent decoder for the compiler's
  own programs was started and then paused, so nothing here establishes portability (MM~\mm{25.209}).
\item Each chapter states what is not known next to the rule it qualifies.
  \Cref{ch:not-known-index} indexes them.
\end{itemize}

\section{Conventions}\label{sec:read-this-document}

Each claim is marked with one of five kinds of statement:

\begin{itemize}
\item \emph{measured}: run on this GPU and compared with an independent reference, with a control that could have failed;
\item \emph{Apple's compiler emits}: read from Apple's compiled code; it shows a pattern works, not that the hardware
  requires it;
\item \emph{decoder}: what Apple's decoder prints or admits, which is not evidence of execution;
\item \emph{this compiler}: a choice the project's compiler makes, not a hardware rule;
\item \emph{open}: not established either way.
\end{itemize}

Opcodes are named first and numbered second, as in "\op{12674}". The numbers are
Apple's opcode ids. The names are this project's, chosen by measured function, and are not Apple's spellings
(\cref{app:c}). Registers are written by hardware number (R0–R125), with their role stated where it matters.

\Cref{tab:reading-guide} lists where to start for each topic.

\begin{xltabular}{\linewidth}{@{}>{\hsize=0.88\hsize}X>{\hsize=1.12\hsize}X@{}}
\caption{Where to start reading, by topic.}\label{tab:reading-guide}\\
\toprule
To understand\dots{} & Read \\
\midrule
\endfirsthead
\tablecontinued{2}
\toprule
To understand\dots{} & Read \\
\midrule
\endhead
how an instruction's bytes become operands & \cref{ch:instruction-encoding}, then \cref{sec:scoreboard}
for the wait fields \\*
register lifetimes, and what a released register returns & \crefrange{sec:operand-modifiers}{sec:loops-loop-variable-rule} \\*
the tensor unit, from the math to the bytes & \cref{ch:tensor-unit}, starting at \cref{sec:fragment-layouts} \\*
why a kernel is slow & \cref{ch:performance} \\*
what orders what in memory & \cref{ch:synchronization}, with \cref{ch:memory} for the
forms \\*
the Metal bypass & \cref{sec:project-authors-two}, then 
\cref{ch:submission-below-metal}, then
\cref{ch:below-metal-native} \\*
how the compiler uses all this & \hyperref[part:ii]{Part II} \\*
a specific opcode & \cref{app:c}, or \cref{app:b} for scalar semantics \\*
\bottomrule
\end{xltabular}
